\par
\SourceLeaf{490}%
I most respectfully leave it to Your — — to consider these reasons and decide on them. I add only that this absence would admittedly put me under some constraint with respect to the two courses I am to give this summer. Given the small number of listeners, however, I could make an arrangement with them to make up the hours that would be missed, partly beforehand and partly afterward. To keep my absence as short as possible, I would make a preliminary verbal agreement with Prof. \textsc{Schumacher} (whom I expect back here any day), so that he could make all the preparations in Lauenburg beforehand and the observations could begin immediately upon my arrival. As the weather at this time of year is not usually unfavorable for observation, I would expect that my absence need last no longer than, say, the trips to spas that some of my colleagues are likely to make around the same time.

I take this opportunity to inform Your — — that the apparatus made by the instrument maker \textsc{Rumpf} for reversing \textsc{Reichenbach}’s transit telescope is now to be mounted in the next few days, and that regular observations with this instrument will then also begin immediately.
\par\nopagebreak[4]
— — —
\par\nopagebreak[4]
Göttingen,\hspace{3em}1819.
\begin{flushright}C. F. \textsc{Gauss}.\end{flushright}

\par\bigskip
\pdfbookmark[0]{Notice on the continuation of the Danish degree surveys. 1820.}{anzeige-gradmessung-1820}
\begin{center}\large NOTICE BY COURT COUNCILLOR GAUSS ON THE CONTINUATION OF THE DANISH DEGREE SURVEYS THROUGH THE KINGDOM OF HANOVER.\end{center}
\par
— — —

Until now, zenith sectors were regarded as an essential requirement even for the foremost observatories. The newly designed meridian circles, particularly \textsc{Reichenbach}’s, now make those instruments dispensable to some extent. At least, in my experience so far, our meridian circle here, in use since February of this year, does everything a zenith sector could do just as well. It might therefore seem unnecessary to give any more thought to a zenith sector for the degree survey, were there not two important circumstances still to be taken into account.

1) The meridian circle requires a very solid mounting, which it can be given only in a proper observatory, and it is to be regarded as a fixed instrument rather than a portable one. The \textit{most important} observations for which a zenith sector or the meridian circle will be used in connection with the degree survey will be those to be made at this observatory, which must in that respect be regarded as the southern endpoint. And observations of this kind on our side at the northern end are, or may seem, superfluous, since those made by the Danish side at Lauenburg can already serve the purpose. Even so, it might prove necessary as the work proceeds, and in any case it will be of scientific interest, to make similar observations at an intermediate point in the Kingdom of Hanover, for example in Celle or Hanover. This could be done only with the aid of a portable zenith sector.

2) Since the use made of the observations at issue here in a degree survey depends entirely on \textit{comparing} the observations made at the various principal points \textit{with one another}, and the observations at Göttingen and perhaps those at Celle or elsewhere on our side will therefore have to be compared with those made by the Danish side at Lauenburg, Lysabbel (on the island of Alsen) and
